(1) 令正传播相位为 \(e^{i(k_xx-\omega t)}\)，上侧场依 z 取 \(e^{-\kappa_dz}\)，下侧场取 \(e^{\kappa_mz}\)，要求 \(\kappa_d,\kappa_m>0\)。Helmholtz 关系分别为
\[
\kappa_d^2=k_x^2-\epsilon_dk_0^2,\qquad
\kappa_m^2=k_x^2-\epsilon_mk_0^2.
\]

(2) TM 模的 H 沿 y，E 在 xz 平面。由 \(\nabla\times H=-i\omega\epsilon_0\epsilon E\) 得两侧切向场幅之间的关系，令 \(H_y,E_x\) 在界面连续，获得
\[
\frac{\kappa_d}{\epsilon_d}+\frac{\kappa_m}{\epsilon_m}=0.
\]
将两侧 κ 关系代入并消元，得到色散
\[
k_x=k_0\sqrt{\frac{\epsilon_d\epsilon_m}{\epsilon_d+\epsilon_m}}.
\]
这是理想无损平界面 TM 表面等离激元的波数。

(3) 需要 kx 实且大于 \(k_0\sqrt{\epsilon_d}\)，保证介质侧指数衰减。由上式得到必要且充分的条件
\[
\epsilon_m<-\epsilon_d.
\]
此时金属侧 \(\kappa_m>0\) 亦自动满足。如果 \(-\epsilon_d<\epsilon_m<0\)，不能同时有实传播波数及两侧束缚衰减。

(4) 当 \(\epsilon_m\to-\infty\)，\(k_x\to k_0\sqrt{\epsilon_d}\)，表面模接近介质光线且介质侧穿透深度增大。当 \(\epsilon_m\to-\epsilon_d^-\)，分母趋零，波数在无损局域模型中形式发散。真实金属的吸收、非局域性与电子量子效应会限制这一无限增大趋势，故不可将无损公式用于无限小的空间尺度。
